\documentclass[11pt,a4paper]{article}
\usepackage[T1]{fontenc}
\usepackage[utf8]{inputenc}
\usepackage{amsmath,amssymb,amsthm}
\usepackage[margin=1in]{geometry}
\usepackage[colorlinks=true,linkcolor=blue,urlcolor=blue]{hyperref}
\theoremstyle{plain}
\newtheorem{theorem}{Theorem}
\newtheorem{lemma}[theorem]{Lemma}
\newtheorem{proposition}[theorem]{Proposition}
\newtheorem{corollary}[theorem]{Corollary}
\theoremstyle{definition}
\newtheorem{remark}[theorem]{Remark}
% A quoted conjecture can be a single hundred-term formula whose terms are thirty-digit
% numbers. TeX cannot hyphenate those, so without a little stretch every such quote runs off
% the page; the linked-a-file papers are all of that shape.
\setlength{\emergencystretch}{3em}

\title{The empirical closed form for OEIS A223930, proved}
\author{Adrian Perez Fontelles\\ \small Independent researcher}
\date{1 September 2026}
\begin{document}
\maketitle

\begin{abstract}
OEIS A223930 carries an empirical closed form for $a(n)$ --- not a recurrence relating
consecutive terms but an explicit formula. It is true, and it is decidable rather than
empirical. The entry counts arrays subject to a condition on a bounded window of consecutive lines, so the lines of an array are the
vertices of a finite digraph, an array is a walk in it, and $a(n)$ is a walk count on
$S=88064$ vertices. Any function of the form $\sum_j p_j(n)\lambda_j^{\,n}$ satisfies the
monic linear recurrence whose characteristic polynomial is
$\prod_j(x-\lambda_j)^{\deg p_j+1}$; here that polynomial is $q(x)=(x-1)^{21}$, of degree
$21$. Two things are then checked exactly: that the walk count satisfies the same
recurrence from some index on, which the Cayley--Hamilton theorem reduces to a finite
computation, and that the two sequences agree at $21$ consecutive indices past that
point. A monic recurrence determines every later term from its predecessors, so agreement
there is agreement for good.
\end{abstract}

\noindent\small 2020 Mathematics Subject Classification. 05A15, 11B37, 15A18.\normalsize

\section{The conjecture}

OEIS A223930 is ``Number of nX5 0..2 arrays with rows, columns and antidiagonals unimodal and diagonals nondecreasing.''. It has offset $1$ and begins
\[
86, 1915, 15791, 86439, 386495, 1548633, 5773556, 20277077,\ \dots
\]
The entry states, as a conjecture contributed by its author and never marked settled:
\begin{quote}\raggedright\small
Empirical: a(n) = (1/1379196149760000)*n\^{}20 - (1/137919614976000)*n\^{}19 + (331/304874938368000)*n\^{}18 - (23/25406244864000)*n\^{}17 + (22273/34871316480000)*n\^{}16 + (2693/697426329600)*n\^{}15 + (709123/2324754432000)*n\^{}14 + (25469/53374464000)*n\^{}13 + (818015239/6897623040000)*n\^{}12 - (447002537/229920768000)*n\^{}11 + (49691417869/1072963584000)*n\^{}10 - (162095712841/268240896000)*n\^{}9 + (48214267849049/6538371840000)*n\^{}8 - (12967593924557/186810624000)*n\^{}7 + (659685543839627/1120863744000)*n\^{}6 - (379579885964341/93405312000)*n\^{}5 + (71087813774265133/3087564480000)*n\^{}4 - (15162068319635473/154378224000)*n\^{}3 + (1768745631701119/6110804700)*n\^{}2 - (6674903392214/14549535)*n + 130847 for n>9
\end{quote}
As of the ``Last modified'' line on the live entry (Oct 03 2025 22:58:31, revision 8) this is still
recorded as empirical, and nothing on the entry records it as proved. Write
\begin{multline*}
f(n)\;=\;\frac{n^{20}}{1379196149760000} - \frac{n^{19}}{137919614976000} \\
+ \frac{331 n^{18}}{304874938368000} - \frac{23 n^{17}}{25406244864000} \\
+ \frac{22273 n^{16}}{34871316480000} + \frac{2693 n^{15}}{697426329600} \\
+ \frac{709123 n^{14}}{2324754432000} + \frac{25469 n^{13}}{53374464000} \\
+ \frac{818015239 n^{12}}{6897623040000} - \frac{447002537 n^{11}}{229920768000} \\
+ \frac{49691417869 n^{10}}{1072963584000} - \frac{162095712841 n^{9}}{268240896000} \\
+ \frac{48214267849049 n^{8}}{6538371840000} - \frac{12967593924557 n^{7}}{186810624000} \\
+ \frac{659685543839627 n^{6}}{1120863744000} - \frac{379579885964341 n^{5}}{93405312000} \\
+ \frac{71087813774265133 n^{4}}{3087564480000} - \frac{15162068319635473 n^{3}}{154378224000} \\
+ \frac{1768745631701119 n^{2}}{6110804700} - \frac{6674903392214 n}{14549535} \\
+ 130847.
\end{multline*}

\section{The count is a walk count}

The entry's defining condition is local: it is a condition on a bounded window of consecutive lines. Take as vertices the
admissible configurations of such a window and put an edge from $u$ to $v$ when $v$ may
follow $u$, which the condition decides by inspection. An admissible array is then exactly a
walk, so with $M$ the adjacency matrix of that digraph on $S=88064$ vertices, and $u,w$ the
vectors recording which windows may start and end an array,
\[
a(n)\;=\;u^{\!\top}M^{\,g(n)}w
\]
for the linear function $g$ that the entry's shape supplies. In particular $a$ is
$C$-finite: it satisfies some monic integer linear recurrence, though not necessarily a short
one.

\section{A closed form is a recurrence}

\begin{lemma}\label{lem:ann}
Let $f(m)=\sum_{j}p_j(m)\lambda_j^{\,m}$ with the $\lambda_j$ distinct and the $p_j$
polynomials, and put $q(x)=\prod_j(x-\lambda_j)^{\deg p_j+1}$. Then $q$ applied to the shift
operator annihilates $f$: writing $q(x)=x^{r}-\sum_{i=1}^{r}c_ix^{r-i}$,
\[
f(m)\;=\;\sum_{i=1}^{r}c_i\,f(m-i)\qquad\text{for every }m.
\]
\end{lemma}

\begin{proof}
The shift operator $E$ acts on $m\mapsto\lambda^{m}p(m)$ as
$(E-\lambda)\bigl(\lambda^{m}p(m)\bigr)=\lambda^{m+1}\bigl(p(m+1)-p(m)\bigr)$, and
$p(m+1)-p(m)$ has degree one less than $p$. So $(E-\lambda)^{\deg p+1}$ kills
$\lambda^{m}p(m)$, and the factors of $q$ commute.
\end{proof}

Here $q(x)=(x-1)^{21}$, of degree $21$; the closed form is a polynomial in $n$, so this is $(x-1)^{21}$, and the recurrence it encodes is
\begin{equation}\label{eq:rec}
a(n)\;=\;21\,a(n-1) - 210\,a(n-2) + 1330\,a(n-3) - 5985\,a(n-4) + 20349\,a(n-5) - 54264\,a(n-6) + 116280\,a(n-7) - 203490\,a(n-8) + 293930\,a(n-9) - 352716\,a(n-10) + 352716\,a(n-11) - 293930\,a(n-12) + 203490\,a(n-13) - 116280\,a(n-14) + 54264\,a(n-15) - 20349\,a(n-16) + 5985\,a(n-17) - 1330\,a(n-18) + 210\,a(n-19) - 21\,a(n-20) + a(n-21).
\end{equation}

\section{The proof}

\begin{lemma}\label{lem:crit}
With $q$ as above, $a$ satisfies \eqref{eq:rec} for all large $n$ if and only if
$u^{\!\top}M^{\,j}q(M)w=0$ for every $j\ge0$, and it is enough to check $j<S$.
\end{lemma}

\begin{proof}
Substituting the walk expression, the residual $a(n)-\sum_ic_ia(n-i)$ equals
$u^{\!\top}M^{\,g(n)-r}q(M)w$ up to the fixed linear reparametrisation $g$. For the bound,
$M^{S}$ is an integer combination of $I,M,\dots,M^{S-1}$ by Cayley--Hamilton, so the values
$u^{\!\top}M^{j}q(M)w$ satisfy the monic recurrence given by the characteristic polynomial of
$M$; $S$ consecutive zeros therefore force all later ones, and the last nonzero value pins the
threshold exactly.
\end{proof}

\begin{theorem}
$a(n)=f(n)$ for every $n\ge10$.
\end{theorem}

\begin{proof}
The vector $w'=q(M)w$ was formed by $21$ matrix--vector products in exact integer
arithmetic, and $u^{\!\top}M^{j}w'$ evaluated for $j=0,1,\dots$, again exactly. Those integers
vanish from the point corresponding to $n=31$ onwards, and the run was continued until the
working vector was identically zero, which settles every larger $j$ at once. So $a$ satisfies
\eqref{eq:rec} for every $n>30$. By Lemma~\ref{lem:ann} $f$ satisfies \eqref{eq:rec}
identically. Both were then evaluated in exact arithmetic at the $21$ consecutive indices
$n=31,\dots,51$ and agree at each. A monic recurrence of order $21$
determines $a(m)$ from $a(m-1),\dots,a(m-21)$, and for every $m\ge52$ both
$m>30$ and $m-21\ge31$ hold, so induction on $m$ gives $a(m)=f(m)$ for every
$m\ge31$.
Below $n=31$ the recurrence argument gives nothing, since \eqref{eq:rec} is only
established for $a$ from $n=31$ on. The remaining indices $n=10,\dots,30$ were
therefore checked one at a time, directly against the walk count in exact integer arithmetic,
and agree.
\end{proof}

The range is tight in the sense that was checked: at $n=9$ the closed form and the sequence differ, so no larger range is available.

\section{Verification}

Four checks, each independent of the algebra above.

First, the walk model reproduces every one of the $21$ terms the entry publishes,
exactly, in integer arithmetic. This is what ties the model to the entry: the model is built
by reading the entry's English, and a misreading gives a different digraph and different
counts.

Second, the closed form was evaluated in exact rational arithmetic --- not floating point ---
at every published term from $n=10$ on, and agrees with the entry's own DATA at each.

Third, the annihilation test of Lemma~\ref{lem:crit} was carried out exactly, and the model
and the closed form were then compared directly at sixty consecutive indices beyond
$n=10$, far past where the proof already settles the matter.

Fourth, the closed form was deliberately perturbed --- by a constant and by a term linear in
$n$ --- and each perturbed form was rejected by the same comparison, so the test is not one
that anything would pass.

\begin{thebibliography}{9}
\bibitem{oeis} The OEIS Foundation, \emph{The On-Line Encyclopedia of Integer
Sequences}, \url{https://oeis.org}, 2026. Sequence A223930.
\bibitem{stanley} R.~P.~Stanley, \emph{Enumerative Combinatorics}, Volume 1, 2nd ed.,
Cambridge University Press, 2011. (Transfer-matrix method, Section 4.7; $C$-finite sequences,
Section 4.1.)
\bibitem{kp} M.~Kauers and P.~Paule, \emph{The Concrete Tetrahedron}, Springer, 2011.
(Closed forms and linear recurrences with constant coefficients.)
\bibitem{horn} R.~A.~Horn and C.~R.~Johnson, \emph{Matrix Analysis}, 2nd ed., Cambridge
University Press, 2013.
\end{thebibliography}

\end{document}
